\documentclass[11pt,letterpaper]{article}

\usepackage[utf8]{inputenc}
\usepackage[T1]{fontenc}
\usepackage[spanish,es-tabla]{babel}
\usepackage{lmodern}
\usepackage{geometry}
\usepackage{xcolor}
\usepackage{graphicx}
\usepackage{amsmath}
\usepackage{booktabs}
\usepackage{longtable}
\usepackage{tabularx}
\usepackage{array}
\usepackage{enumitem}
\usepackage{hyperref}
\usepackage{fancyhdr}
\usepackage{lastpage}
\usepackage{titlesec}
\usepackage{tikz}
\usetikzlibrary{arrows.meta,positioning,shapes.geometric,calc,backgrounds,fit,babel}

\geometry{margin=2.35cm}
\hypersetup{
  colorlinks=true,
  linkcolor=black,
  urlcolor=blue,
  pdftitle={Roadmap SIMIO - Capstone Analytics},
  pdfauthor={Joaquín Parraud, Sebastián Ramírez y Vicente Ramírez}
}

\definecolor{uddblue}{HTML}{003B5C}
\definecolor{uddgray}{HTML}{4D4D4D}
\definecolor{softgray}{HTML}{F2F4F7}
\definecolor{linegray}{HTML}{D9DEE7}
\definecolor{accentgreen}{HTML}{2E7D32}
\definecolor{accentorange}{HTML}{C75B12}
\definecolor{accentred}{HTML}{B3261E}

\tikzset{
  box/.style={draw=uddblue, fill=softgray, rounded corners=2pt, align=center, minimum height=0.78cm, text width=2.45cm, font=\scriptsize},
  smallbox/.style={draw=uddblue, fill=softgray, rounded corners=2pt, align=center, minimum height=0.65cm, text width=2.05cm, font=\scriptsize},
  decision/.style={diamond, draw=accentorange, fill=orange!8, aspect=1.8, align=center, inner sep=1.2pt, font=\scriptsize},
  arrow/.style={-Latex, thick, draw=uddblue},
  dashedarrow/.style={-Latex, thick, dashed, draw=uddgray},
  lane/.style={draw=linegray, fill=softgray, rounded corners=3pt},
  ganttbar/.style={draw=uddblue, fill=uddblue!25, rounded corners=1pt},
  ganttmilestone/.style={draw=accentred, fill=accentred!75}
}

\setlength{\parindent}{0pt}
\setlength{\parskip}{0.45em}
\setlist[itemize]{leftmargin=1.2em,itemsep=0.15em,topsep=0.2em}
\setlist[enumerate]{leftmargin=1.4em,itemsep=0.15em,topsep=0.2em}
\renewcommand{\arraystretch}{1.18}
\newcolumntype{Y}{>{\raggedright\arraybackslash}X}
\newcolumntype{L}[1]{>{\raggedright\arraybackslash}p{#1}}

\titleformat{\section}{\Large\bfseries\color{uddblue}}{\thesection}{0.75em}{}
\titleformat{\subsection}{\large\bfseries\color{uddgray}}{\thesubsection}{0.75em}{}

\pagestyle{fancy}
\fancyhf{}
\lhead{\small Capstone Analytics IIC511AD}
\rhead{\small Roadmap SIMIO}
\cfoot{\small Página \thepage\ de \pageref{LastPage}}

\newcommand{\curso}{Capstone Analytics}
\newcommand{\codigo}{IIC511AD}
\newcommand{\profesor}{Victor Vera}
\newcommand{\universidad}{Universidad del Desarrollo}
\newcommand{\fechaemision}{27 de abril de 2026}
\newcommand{\horizonte}{22 de abril de 2026 al 22 de mayo de 2026}
\newcommand{\versiondoc}{v1.0}

\begin{document}

\begin{titlepage}
  \centering
  \vspace*{0.5cm}
  \fbox{%
    \begin{minipage}[c][2.8cm][c]{0.42\textwidth}
      \centering
      \textcolor{uddgray}{Espacio reservado para logo institucional}
    \end{minipage}
  }

  \vspace{1.2cm}
  {\Large \textbf{\universidad}\par}
  \vspace{0.25cm}
  {\large \curso\ (\codigo)\par}
  {\large Profesor: \profesor\par}

  \vspace{1.4cm}
  {\Huge \textbf{Proyecto Capstone Analytics}\par}
  \vspace{0.45cm}
  {\LARGE \textbf{Roadmap profesional de simulación en SIMIO}\par}
  \vspace{0.3cm}
  {\Large Simulación de Eventos Discretos: Proceso de Fabricación de Pan\par}

  \vspace{1.0cm}
  \begin{tabular}{rl}
    \textbf{Horizonte de ejecución:} & \horizonte \\
    \textbf{Fecha de emisión:} & \fechaemision \\
    \textbf{Versión:} & \versiondoc \\
    \textbf{Rama de trabajo:} & \texttt{Simu/Origin} \\
  \end{tabular}

  \vfill
  \begin{tabular}{ll}
    \toprule
    \textbf{Integrante} & \textbf{Responsabilidad principal} \\
    \midrule
    Joaquín Parraud & Modelo físico, layout operacional y construcción en SIMIO \\
    Sebastián Ramírez & Flujo lógico, reglas de proceso y verificación estructural \\
    Vicente Ramírez & Análisis de datos, demanda y triangulación de probabilidades \\
    \bottomrule
  \end{tabular}

  \vspace*{0.5cm}
\end{titlepage}

\tableofcontents
\newpage

\section{Control del Documento}

\begin{tabularx}{\textwidth}{L{2.2cm}L{3.0cm}L{4.0cm}Y}
  \toprule
  \textbf{Versión} & \textbf{Fecha} & \textbf{Responsable} & \textbf{Descripción} \\
  \midrule
  \versiondoc & \fechaemision & Equipo de proyecto & Reestructuración profesional del roadmap, incorporación de gobernanza, trazabilidad, validación, matriz experimental y responsabilidades nominales. \\
  \bottomrule
\end{tabularx}

\textbf{Estado del documento:} planificación ejecutiva y técnica para implementación, validación y defensa del modelo de simulación.

\textbf{Propósito:} establecer una ruta de trabajo verificable para construir un modelo SIMIO ejecutable, estadísticamente defendible y útil para recomendar capacidad, dotación y política operacional de una panadería de supermercado.

\section{Resumen Ejecutivo}

El proyecto busca modelar, mediante simulación de eventos discretos en SIMIO, la operación de una panadería de supermercado que produce diez tipos de pan con recursos compartidos, demanda horaria no uniforme y riesgo de quiebre de stock. El objetivo no es únicamente reproducir el proceso, sino entregar una recomendación operacional defendible: cuántos recursos utilizar, qué política de horneo adoptar y cómo balancear servicio al cliente con costo operacional.

La propuesta se estructura en cuatro ejes: modelación física, lógica de proceso, análisis de datos y validación estadística. Joaquín Parraud liderará la construcción física del modelo en SIMIO; Sebastián Ramírez diseñará el flujo lógico, las reglas de producción y las pruebas de consistencia; Vicente Ramírez liderará el análisis de datos, el motor de demanda y la triangulación de probabilidades. Las actividades restantes se distribuyen de manera uniforme para evitar dependencia crítica de una sola persona y asegurar continuidad del trabajo.

El criterio rector será la trazabilidad. Cada decisión relevante del modelo deberá quedar asociada a una fuente de datos, una regla operacional, un indicador de desempeño y una evidencia verificable. De este modo, la entrega final podrá sostenerse tanto ante una audiencia técnica como ante una audiencia gerencial.

\section{Mandato, Alcance y Criterios de Éxito}

\subsection{Mandato del Proyecto}

El sistema representa una panadería que abastece una sala de ventas entre las 09:00 y las 21:00 horas. Antes de la apertura debe existir stock inicial de todos los productos, y durante el día se debe producir, hornear, enfriar y reponer pan para reducir quiebres. Si un cliente no encuentra el pan requerido, la venta se pierde y no existe sustitución.

El estudio debe determinar:
\begin{itemize}
  \item dotación recomendada de panaderos y manipuladores;
  \item número requerido de hornos, mezcladoras, amasadoras, estaciones, cámaras, carros y bandejas;
  \item horario operativo conveniente de la panadería;
  \item política de secuencia de producción y horneo;
  \item nivel de servicio alcanzable bajo costo operacional controlado.
\end{itemize}

\subsection{Alcance Mínimo del Modelo}

El modelo debe representar producción, recursos limitados, inventario en sala, turnos, descansos, demanda no uniforme, compatibilidad de familias de horno, setup por cambio de familia, enfriamiento, reposición y pérdida de venta por quiebre. Además, debe permitir comparar políticas alternativas de producción y secuencia de horno.

\subsection{Criterios de Éxito}

\begin{tabularx}{\textwidth}{L{4.1cm}Y}
  \toprule
  \textbf{Criterio} & \textbf{Evidencia esperada} \\
  \midrule
  Ejecutabilidad & Modelo SIMIO corre sin errores y reproduce el escenario base. \\
  Coherencia operacional & No hay mezclas de familias incompatibles, retiros parciales de horno ni inventarios negativos. \\
  Calidad estadística & Resultados con réplicas, media, desviación estándar e intervalos de confianza. \\
  Trazabilidad & Cada KPI se vincula a fuente de datos, regla de modelo y salida de resultados. \\
  Decisión gerencial & Recomendación final presenta trade-off explícito entre costo, servicio y riesgo. \\
  \bottomrule
\end{tabularx}

\section{Equipo, Roles y Gobernanza}

\subsection{Responsabilidades Principales}

\begin{tabularx}{\textwidth}{L{3.3cm}L{4.4cm}Y}
  \toprule
  \textbf{Integrante} & \textbf{Rol principal} & \textbf{Responsabilidades complementarias} \\
  \midrule
  Joaquín Parraud & Modelo físico en SIMIO & Construcción de layout, recursos, estaciones, caminos, objetos, capacidades, animación operacional y apoyo en pruebas de corridas. \\
  Sebastián Ramírez & Flujo lógico del proceso & Secuencia productiva, reglas de inventario, política de hornos, restricciones operativas, validación estructural y documentación de lógica. \\
  Vicente Ramírez & Datos y triangulación de probabilidades & Limpieza de datos, demanda horaria, probabilidades de elección, calibración, diseño experimental, análisis estadístico y visualización de resultados. \\
  \bottomrule
\end{tabularx}

\subsection{RACI de Macroentregables}

\begin{tabularx}{\textwidth}{L{4.8cm}L{2.4cm}L{2.4cm}L{2.4cm}Y}
  \toprule
  \textbf{Entregable} & \textbf{Joaquín} & \textbf{Sebastián} & \textbf{Vicente} & \textbf{Profesor} \\
  \midrule
  Arquitectura física SIMIO & R/A & C & C & I \\
  Flujo lógico y reglas operativas & C & R/A & C & I \\
  Motor de demanda y probabilidades & C & C & R/A & I \\
  Política de horno y secuenciación & C & R/A & C & I \\
  Diseño experimental y análisis estadístico & C & C & R/A & I \\
  Verificación y validación & R & R/A & R & C \\
  Informe técnico gerencial & C & R & R/A & I \\
  Informe académico y anexos & R & R & R/A & I \\
  Recomendación final costo-servicio & A & A & A & C \\
  \bottomrule
\end{tabularx}

\textit{Leyenda: R = Responsible; A = Accountable; C = Consulted; I = Informed.}

\section{Arquitectura de Datos y Trazabilidad}

\subsection{Fuentes de Entrada}

\begin{tabularx}{\textwidth}{L{4.0cm}L{4.6cm}Y}
  \toprule
  \textbf{Archivo} & \textbf{Contenido} & \textbf{Uso en el modelo} \\
  \midrule
  \parbox[t]{4.0cm}{\texttt{parametros\_proceso}\\\texttt{\_panaderia.xlsx}} & Tiempos de etapas, familias, lotes, temperaturas y capacidad por carro. & Parametrizar procesos, tiempos, familias de horno, tamaños de lote y requerimientos de capacidad. \\
  \parbox[t]{4.0cm}{\texttt{perfil\_demanda}\\\texttt{\_por\_hora}\\\texttt{\_panaderia.xlsx}} & Kilogramos esperados por hora y tipo de pan. & Calibrar intensidad de demanda, validación por franja y volumen diario. \\
  \parbox[t]{4.0cm}{\texttt{probabilidades}\\\texttt{\_eleccion}\\\texttt{\_por\_hora.xlsx}} & Probabilidad horaria de elección de cada producto. & Simular composición de compra, primera elección y elecciones adicionales sin repetición. \\
  Enunciado del caso & Reglas, restricciones, entregables e indicadores. & Definir alcance mínimo, supuestos obligatorios y criterios de aceptación. \\
  \bottomrule
\end{tabularx}

\subsection{Matriz de Trazabilidad}

\small
\begin{longtable}{L{2.55cm}L{2.75cm}L{3.05cm}L{2.75cm}L{2.75cm}}
  \toprule
  \textbf{Requisito} & \textbf{Fuente} & \textbf{Implementación SIMIO} & \textbf{KPI asociado} & \textbf{Evidencia} \\
  \midrule
  \endfirsthead
  \toprule
  \textbf{Requisito} & \textbf{Fuente} & \textbf{Implementación SIMIO} & \textbf{KPI asociado} & \textbf{Evidencia} \\
  \midrule
  \endhead
  Inventario por tipo de pan & Enunciado, sección C & Estados por producto en sala & Fill rate, quiebre, sobrante & Tabla de inventario horario \\
  Demanda horaria no uniforme & Perfil de demanda & Llegadas/órdenes por franja & Demanda simulada vs. esperada & Validación con tolerancia \(\pm 5\%\) \\
  Compra de 1, 2 o 3 tipos & Enunciado, sección E & Lógica probabilística sin repetición & Mix de compra y kg vendidos & Log de demanda y elecciones \\
  Familias de horneado & Enunciado, sección III-B & Atributo \texttt{FamiliaHorno} & Utilización y setups & Historial de corridas \\
  No mezclar familias & Enunciado, sección III-C & Restricción de carga por familia & Violaciones lógicas & Prueba de compatibilidad \\
  Carga/descarga con manipulador & Enunciado, sección III-E & Recurso requerido antes y después del horno & Espera de horno, utilización manipuladores & Reporte de colas \\
  Setup por cambio de familia & Enunciado, sección III-F & Tiempo condicional de 0 o 5 min & Tiempo ocioso, throughput & Registro de secuencias \\
  Turnos y descansos escalonados & Enunciado, sección D & Calendarios por recurso humano & Disponibilidad y espera & Calendario auditado \\
  Enfriamiento no disponible para venta & Enunciado, sección III-J & Estado intermedio antes de inventario sala & Tiempo de ciclo y fill rate & Tiempos por lote \\
  Recursos limitados & Enunciado, sección B & Capacidades parametrizadas por escenario & Utilización y colas & Comparativo DoE \\
  \bottomrule
\end{longtable}
\normalsize

\section{Arquitectura Operacional del Modelo}

\subsection{Entidad Principal}

La entidad \texttt{LotePan} representará una unidad de producción. Sus atributos mínimos serán: tipo de pan, familia de horneado, kilogramos del lote, hora de liberación, hora de llegada a sala, prioridad operacional, tiempo acumulado de proceso y estado de disponibilidad.

\subsection{Flujo Lógico de Proceso}

El flujo lógico será:

\begin{enumerate}
  \item solicitud de producción por política de inventario y demanda esperada;
  \item pesado y dosificación;
  \item mezclado o amasado según corresponda;
  \item reposo en masa;
  \item dividido, formado o laminado;
  \item fermentación final;
  \item espera de horno por familia compatible;
  \item setup si cambia la familia;
  \item carga, horneado y descarga;
  \item enfriamiento;
  \item reposición a sala;
  \item consumo por demanda y registro de venta satisfecha o perdida.
\end{enumerate}

\begin{figure}[htbp]
  \centering
  \resizebox{\textwidth}{!}{%
  \begin{tikzpicture}[node distance=0.7cm and 0.55cm]
    \node[smallbox] (pesado) {Pesado y\\dosificación};
    \node[smallbox, right=of pesado] (mezclado) {Mezclado};
    \node[smallbox, right=of mezclado] (amasado) {Amasado};
    \node[smallbox, right=of amasado] (reposo) {Reposo\\en masa};
    \node[smallbox, right=of reposo] (formado) {Dividido y\\formado};

    \node[smallbox, below=1.0cm of formado] (fermentacion) {Fermentación\\final};
    \node[smallbox, left=of fermentacion] (horno) {Horneado\\batch};
    \node[smallbox, left=of horno] (enfriado) {Enfriado};
    \node[smallbox, left=of enfriado] (reposicion) {Reposición\\a sala};
    \node[smallbox, left=of reposicion] (inventario) {Inventario\\en sala};

    \node[decision, below=0.85cm of inventario] (stock) {Stock\\disponible?};
    \node[smallbox, below left=0.75cm and 0.05cm of stock, fill=accentred!8, draw=accentred] (quiebre) {Registrar\\quiebre};
    \node[smallbox, below right=0.75cm and 0.05cm of stock, fill=accentgreen!8, draw=accentgreen] (venta) {Satisfacer\\demanda};

    \draw[arrow] (pesado) -- (mezclado);
    \draw[arrow] (mezclado) -- (amasado);
    \draw[arrow] (amasado) -- (reposo);
    \draw[arrow] (reposo) -- (formado);
    \draw[arrow] (formado) -- (fermentacion);
    \draw[arrow] (fermentacion) -- (horno);
    \draw[arrow] (horno) -- (enfriado);
    \draw[arrow] (enfriado) -- (reposicion);
    \draw[arrow] (reposicion) -- (inventario);
    \draw[arrow] (inventario) -- (stock);
    \draw[arrow] (stock) -- node[font=\scriptsize, left] {No} (quiebre);
    \draw[arrow] (stock) -- node[font=\scriptsize, right] {Sí} (venta);
  \end{tikzpicture}}
  \caption{Flujo lógico end-to-end desde preparación hasta venta o quiebre.}
\end{figure}

\begin{figure}[htbp]
  \centering
  \resizebox{\textwidth}{!}{%
  \begin{tikzpicture}[node distance=0.55cm and 0.65cm]
    \node[box] (panaderos) {Panaderos};
    \node[box, right=of panaderos] (prep) {Cola de\\preparación};
    \node[box, right=of prep] (mezcla) {Mezcladoras\\y amasadoras};
    \node[box, right=of mezcla] (mesas) {Mesas de\\formado};

    \node[box, below=1.05cm of mesas] (camaras) {Cámaras de\\fermentación};
    \node[box, left=of camaras] (manip) {Manipuladores};
    \node[box, left=of manip] (colahorno) {Cola de\\horneado};
    \node[box, left=of colahorno] (hornos) {Hornos\\+ setup};

    \node[box, below=1.05cm of hornos] (enfriar) {Cola de\\enfriado};
    \node[box, right=of enfriar] (racks) {Carros y\\bandejas};
    \node[box, right=of racks] (sala) {Inventario\\sala};
    \node[box, right=of sala] (clientes) {Demanda\\clientes};

    \draw[arrow] (panaderos) -- (prep);
    \draw[arrow] (prep) -- (mezcla);
    \draw[arrow] (mezcla) -- (mesas);
    \draw[arrow] (mesas) -- (camaras);
    \draw[arrow] (camaras) -- (colahorno);
    \draw[arrow] (colahorno) -- (hornos);
    \draw[arrow] (hornos) -- (enfriar);
    \draw[arrow] (enfriar) -- (racks);
    \draw[arrow] (racks) -- (sala);
    \draw[arrow] (sala) -- (clientes);
    \draw[dashedarrow] (manip) -- (hornos);
    \draw[dashedarrow] (manip) -- (racks);

    \begin{scope}[on background layer]
      \node[lane, fit=(panaderos)(prep)(mezcla)(mesas), label={[font=\small\bfseries]above:Producción de masa y formado}] {};
      \node[lane, fit=(camaras)(manip)(colahorno)(hornos), label={[font=\small\bfseries]above:Fermentación y horneo}] {};
      \node[lane, fit=(enfriar)(racks)(sala)(clientes), label={[font=\small\bfseries]below:Salida a ventas}] {};
    \end{scope}
  \end{tikzpicture}}
  \caption{Arquitectura de recursos, colas y dependencias operacionales.}
\end{figure}

\subsection{Política de Horno}

La política base seguirá la regla de carga mínima recomendada, espera máxima y riesgo de quiebre:

\begin{enumerate}
  \item agrupar lotes disponibles por familia compatible;
  \item evaluar carga acumulada frente al umbral mínimo;
  \item si la carga alcanza el umbral, iniciar corrida;
  \item si no alcanza el umbral, esperar acumulación hasta el máximo permitido;
  \item si existe riesgo de quiebre, autorizar corrida anticipada;
  \item si no existe urgencia, postergar y reevaluar secuencia.
\end{enumerate}

\begin{figure}[htbp]
  \centering
  \resizebox{0.92\textwidth}{!}{%
  \begin{tikzpicture}[node distance=0.85cm and 1.25cm]
    \node[box] (solicitud) {Solicitud de\\corrida};
    \node[decision, below=of solicitud] (carga) {Carga\\\(\geq\) umbral?};
    \node[box, right=1.35cm of carga, fill=accentgreen!8, draw=accentgreen] (iniciar) {Iniciar corrida\\familia prioritaria};
    \node[decision, below=of carga] (espera) {Espera\\\(<\) máximo?};
    \node[box, right=1.35cm of espera] (acumular) {Esperar\\acumulación};
    \node[decision, below=of espera] (riesgo) {Riesgo de\\quiebre alto?};
    \node[box, right=1.35cm of riesgo, fill=accentorange!8, draw=accentorange] (anticipar) {Lanzar corrida\\anticipada};
    \node[box, below=of riesgo] (postergar) {Postergar y\\reordenar};

    \draw[arrow] (solicitud) -- (carga);
    \draw[arrow] (carga) -- node[font=\scriptsize, above] {Sí} (iniciar);
    \draw[arrow] (carga) -- node[font=\scriptsize, left] {No} (espera);
    \draw[arrow] (espera) -- node[font=\scriptsize, above] {Sí} (acumular);
    \draw[arrow] (espera) -- node[font=\scriptsize, left] {No} (riesgo);
    \draw[arrow] (riesgo) -- node[font=\scriptsize, above] {Sí} (anticipar);
    \draw[arrow] (riesgo) -- node[font=\scriptsize, left] {No} (postergar);
    \draw[dashedarrow] (acumular.north) .. controls +(0,1.0) and +(1.0,-0.5) .. (carga.east);
  \end{tikzpicture}}
  \caption{Árbol de decisión para activación y secuencia de corrida de horno.}
\end{figure}

\textbf{Nota crítica de consistencia:} el enunciado menciona carga mínima recomendada de 50\% equivalente a 600 kg. Sin embargo, el archivo de parámetros contiene capacidades referenciales por carro entre 28 y 45 kg; con 6 carros, la capacidad máxima por corrida estaría aproximadamente entre 168 y 270 kg, dependiendo del producto. Esta diferencia debe resolverse antes de cerrar el modelo final, ya sea mediante consulta al profesor, corrección de unidades o supuesto explícito documentado.

\section{Motor de Demanda y Triangulación de Probabilidades}

Vicente Ramírez liderará el diseño y validación del motor de demanda. La lógica propuesta será:

\begin{enumerate}
  \item para cada franja horaria, definir el volumen esperado por tipo de pan;
  \item generar compras de clientes con probabilidad de adquirir 1, 2 o 3 tipos;
  \item seleccionar el primer tipo de pan según probabilidades horarias;
  \item seleccionar tipos adicionales con las mismas proporciones, excluyendo productos ya escogidos;
  \item generar kilogramos por tipo mediante distribución triangular parametrizada por producto;
  \item comparar la demanda simulada acumulada contra el perfil esperado por hora;
  \item ajustar el número de clientes o intensidad de llegadas hasta respetar la tolerancia definida.
\end{enumerate}

La validación de demanda se considerará aceptable si el error por franja y por producto crítico permanece dentro de una tolerancia operacional de \(\pm 5\%\), salvo justificación estadística documentada.

\section{Indicadores de Desempeño}

\begin{longtable}{L{4.1cm}L{5.2cm}L{3.2cm}L{2.6cm}}
  \toprule
  \textbf{KPI} & \textbf{Definición} & \textbf{Meta o criterio} & \textbf{Uso} \\
  \midrule
  \endfirsthead
  \toprule
  \textbf{KPI} & \textbf{Definición} & \textbf{Meta o criterio} & \textbf{Uso} \\
  \midrule
  \endhead
  Fill rate total & Demanda satisfecha sobre demanda total & \(\geq 95\%\) & Servicio global \\
  Fill rate por tipo & Demanda satisfecha por producto & Prioridad en productos críticos & Servicio detallado \\
  Kg no vendidos por quiebre & Demanda perdida por falta de stock & Minimizar & Impacto comercial \\
  Sobrante fin de día & Kg remanentes al cierre & \(\leq 8\%\) de producción & Control de sobreproducción \\
  Utilización de panaderos & Tiempo ocupado sobre disponibilidad & 65\%--90\% & Balance humano \\
  Utilización de manipuladores & Tiempo ocupado sobre disponibilidad & 65\%--90\% & Riesgo de bloqueo \\
  Utilización de hornos & Tiempo de horno activo sobre disponibilidad & 60\%--92\% & Capacidad de horneo \\
  Tiempo ciclo por lote & Desde liberación hasta disponibilidad en sala & Minimizar con estabilidad & Fluidez del proceso \\
  Espera por recurso & Tiempo promedio en cola por etapa & Minimizar & Cuello de botella \\
  Producción por tipo & Kg fabricados por producto & Coherente con demanda y sobrante & Control productivo \\
  \bottomrule
\end{longtable}

\section{Modelo de Costo Relativo}

Si no se cuenta con costos unitarios reales, se utilizará un índice de costo relativo para comparar escenarios:

\[
\begin{aligned}
ICR ={}& w_1 \cdot \text{Panaderos}
+ w_2 \cdot \text{Manipuladores}
+ w_3 \cdot \text{Hornos} \\
&+ w_4 \cdot \text{Equipos}
+ w_5 \cdot \text{Sobrante}
+ w_6 \cdot \text{Kg no vendidos}
\end{aligned}
\]

Los pesos \(w_i\) deberán declararse en el informe. La recomendación final no se basará solo en menor costo, sino en la frontera costo-servicio: configuraciones que logren alto fill rate, baja pérdida por quiebre y utilización estable sin sobredimensionar recursos.

\section{Plan de Verificación y Validación}

\subsection{Verificación Estructural}

\begin{itemize}
  \item El modelo ejecuta sin errores ni entidades bloqueadas.
  \item No existen inventarios negativos.
  \item No se mezclan familias incompatibles en una corrida de horno.
  \item No hay retiro parcial de carga antes de finalizar horneado y descarga.
  \item Los descansos no dejan una función crítica sin cobertura.
  \item Todo lote vendido pasó por enfriamiento antes de llegar a sala.
\end{itemize}

\subsection{Validación Cuantitativa}

\begin{itemize}
  \item La demanda diaria simulada converge a 8.200 kg/día en el escenario base.
  \item El perfil horario simulado respeta el patrón de los archivos de entrada.
  \item La producción, venta, sobrante y quiebre cierran mediante balance de masa.
  \item Los KPIs se reportan por tipo de pan y total.
  \item Cada escenario se ejecuta con al menos 30 réplicas e intervalos de confianza al 95\%.
\end{itemize}

\section{Diseño Experimental}

\begin{tabularx}{\textwidth}{L{4.2cm}L{4.2cm}Y}
  \toprule
  \textbf{Factor} & \textbf{Niveles sugeridos} & \textbf{Motivo} \\
  \midrule
  Panaderos & 2, 3, 4 & Determina capacidad de preparación, formado y apoyo. \\
  Manipuladores & 1, 2, 3 & Afecta carga, descarga, reposición y bloqueo de horno. \\
  Hornos & 1, 2 & Evalúa capacidad crítica de horneo. \\
  Mezcladoras / amasadoras & 1, 2 & Identifica restricciones tempranas del proceso. \\
  Política de secuencia & Fija, por stock, por demanda, híbrida & Compara estabilidad con respuesta a quiebres. \\
  Carga mínima efectiva & 40\%, 50\%, 60\% & Evalúa eficiencia versus velocidad de respuesta. \\
  Espera máxima de corrida & 10, 15, 20 min & Controla acumulación y urgencia operacional. \\
  \bottomrule
\end{tabularx}

\begin{figure}[htbp]
  \centering
  \resizebox{0.88\textwidth}{!}{%
  \begin{tikzpicture}[x=1cm,y=0.06cm]
    \draw[-Latex, thick] (0,0) -- (0,105) node[above, font=\scriptsize] {Índice normalizado};
    \draw[-Latex, thick] (0,0) -- (10.2,0);
    \foreach \y in {20,40,60,80,100} {
      \draw[linegray] (0,\y) -- (9.7,\y);
      \node[left, font=\scriptsize] at (0,\y) {\y};
    }
    \foreach \x/\lab in {1/Base,3/S1,5/S2,7/S3} {
      \node[below, font=\scriptsize] at (\x+0.45,0) {\lab};
    }
    \foreach \x/\a/\b/\c in {1/82/70/28,3/91/82/16,5/96/88/8,7/94/78/11} {
      \draw[fill=accentgreen!65, draw=accentgreen] (\x,0) rectangle +(0.28,\a);
      \draw[fill=uddblue!55, draw=uddblue] (\x+0.35,0) rectangle +(0.28,\b);
      \draw[fill=accentred!55, draw=accentred] (\x+0.70,0) rectangle +(0.28,\c);
    }
    \draw[fill=accentgreen!65, draw=accentgreen] (8.7,93) rectangle (8.95,97);
    \node[font=\scriptsize, anchor=west] at (9.05,95) {Fill rate};
    \draw[fill=uddblue!55, draw=uddblue] (8.7,83) rectangle (8.95,87);
    \node[font=\scriptsize, anchor=west] at (9.05,85) {Costo};
    \draw[fill=accentred!55, draw=accentred] (8.7,73) rectangle (8.95,77);
    \node[font=\scriptsize, anchor=west] at (9.05,75) {Quiebre};
  \end{tikzpicture}}
  \caption{Comparativo esperado de escenarios para comunicar el trade-off costo-servicio.}
\end{figure}

\textbf{Regla de reporte:} cada escenario tendrá ID, parámetros, semilla, número de réplicas, fecha de corrida, versión del modelo, KPIs principales y observaciones. Esta bitácora será anexo obligatorio.

\section{Roadmap de Ejecución}

\begin{figure}[htbp]
  \centering
  \resizebox{\textwidth}{!}{%
  \begin{tikzpicture}[x=0.34cm,y=0.72cm]
    \draw[-Latex, thick] (0,0) -- (32,0);
    \foreach \x/\lab in {0/22-abr,8/30-abr,9/01-may,16/08-may,23/15-may,29/21-may,30/22-may} {
      \draw (\x,0.12) -- (\x,-0.12);
      \node[below, font=\scriptsize] at (\x,-0.15) {\lab};
    }
    \foreach \y/\name in {1/Gobernanza y diseño,2/Implementación y calibración,3/Experimentación,4/Documentación y cierre} {
      \node[left, align=right, font=\scriptsize] at (0,\y) {\name};
      \draw[linegray] (0,\y) -- (30.5,\y);
    }
    \draw[ganttbar] (0,0.78) rectangle (9,1.22);
    \draw[ganttbar] (9,1.78) rectangle (17,2.22);
    \draw[ganttbar] (17,2.78) rectangle (24,3.22);
    \draw[ganttbar] (24,3.78) rectangle (30,4.22);
    \foreach \x/\m in {9/M1,16/M2,23/M3,29/M4,30/M5} {
      \draw[ganttmilestone] (\x-0.16,0.55) rectangle (\x+0.16,4.45);
      \node[above, font=\scriptsize\bfseries, text=accentred] at (\x,4.55) {\m};
    }
  \end{tikzpicture}}
  \caption{Gantt ejecutivo del proyecto y gates de control.}
\end{figure}

\subsection{Calendario Maestro}

\begin{longtable}{L{2.4cm}L{4.5cm}L{3.6cm}L{4.3cm}}
  \toprule
  \textbf{Fecha} & \textbf{Foco} & \textbf{Responsable líder} & \textbf{Entregable de cierre} \\
  \midrule
  \endfirsthead
  \toprule
  \textbf{Fecha} & \textbf{Foco} & \textbf{Responsable líder} & \textbf{Entregable de cierre} \\
  \midrule
  \endhead
  22-abr & Kickoff técnico, lectura del enunciado y KPIs objetivo & Equipo & Charter operacional y criterios de calidad \\
  23-abr & Mapeo As-Is / To-Be y eventos discretos & Sebastián & Diagrama lógico validado \\
  24-abr & Normalización de datos y diccionario & Vicente & Dataset limpio y diccionario de variables \\
  27-abr & Estructura física base en SIMIO & Joaquín & Layout y recursos base creados \\
  28-abr & Motor de demanda horaria y elección de productos & Vicente & Demanda simulada calibrada preliminar \\
  29-abr & Inventario en sala, quiebres y ventas perdidas & Sebastián & Subsistema de ventas operativo \\
  30-abr & Turnos, descansos y cobertura mínima & Sebastián + Joaquín & Calendarios sin solapes críticos \\
  01-may & Gate M1: verificación estructural & Equipo & Checklist estructural cerrado \\
  04-may & Flujo completo mezcla--venta & Joaquín + Sebastián & Modelo end-to-end ejecutable \\
  05-may & Política de horno, setup y compatibilidad & Sebastián & Motor de corridas operativo \\
  06-may & Recursos limitados: carros, bandejas y equipos & Joaquín & Capacidades parametrizadas \\
  07-may & Calibración base-case 8.200 kg/día & Vicente & Escenario base estable \\
  08-may & Gate M2: validación cuantitativa inicial & Vicente + Sebastián & Informe de brechas y ajustes \\
  11-may & Matriz DoE y escenarios priorizados & Vicente & Catálogo de escenarios \\
  12-may & Corridas masivas con réplicas e IC95\% & Vicente + Joaquín & Resultados comparables \\
  13-may & Sensibilidad y cuello de botella & Equipo & Ranking de restricciones \\
  14-may & Selección de configuraciones candidatas & Equipo & Top 3 costo-servicio \\
  15-may & Gate M3: recomendación preliminar & Equipo & Recomendación defendible \\
  18-may & Informe técnico gerencial & Sebastián + Vicente & Borrador técnico completo \\
  19-may & Informe académico y anexos & Joaquín + Vicente & Borrador académico completo \\
  20-may & Integración final, trazabilidad y bitácora & Equipo & Paquete documental consolidado \\
  21-may & Ensayo de defensa y revisión cruzada & Equipo & Versión corregida \\
  22-may & Gate M4/M5: cierre y entrega & Equipo & Entrega final versionada \\
  \bottomrule
\end{longtable}

\subsection{Balance de Carga por Integrante}

La asignación busca respetar especialización sin concentrar riesgo. Joaquín concentra construcción física y parametrización de recursos; Sebastián concentra lógica, restricciones y verificación; Vicente concentra datos, experimentación y análisis. Las etapas de validación, documentación y defensa se trabajan de forma cruzada para que todos puedan explicar el modelo completo.

\section{Riesgos y Contramedidas}

\begin{longtable}{L{0.9cm}L{4.45cm}L{4.65cm}L{2.0cm}L{1.95cm}}
  \toprule
  \textbf{ID} & \textbf{Riesgo} & \textbf{Contramedida} & \textbf{Owner} & \textbf{Trigger} \\
  \midrule
  \endfirsthead
  \toprule
  \textbf{ID} & \textbf{Riesgo} & \textbf{Contramedida} & \textbf{Owner} & \textbf{Trigger} \\
  \midrule
  \endhead
  R1 & Inconsistencia entre capacidad declarada y archivo de parámetros. & Documentar supuesto y consultar al profesor antes de congelar escenarios. & Sebastián & Diferencia no resuelta al M2 \\
  R2 & Demanda simulada no reproduce perfil horario. & Ajustar intensidad de clientes y validar por franja con tolerancia. & Vicente & Error \(>5\%\) \\
  R3 & Bloqueo de horno por falta de manipulador. & Escalonar turnos y evaluar prioridad de carga/descarga. & Joaquín & Espera horno \(>10\) min \\
  R4 & Sobreproducción por política agresiva. & Penalizar sobrante en índice de costo relativo. & Vicente & Sobrante \(>8\%\) \\
  R5 & Modelo físico avanza sin trazabilidad documental. & Actualizar matriz de requisitos en cada gate. & Sebastián & Cambio sin registro \\
  R6 & Dependencia excesiva de una persona. & Revisión cruzada semanal y explicación breve de cada módulo al equipo. & Equipo & Módulo no explicable por otro integrante \\
  \bottomrule
\end{longtable}

\begin{figure}[htbp]
  \centering
  \resizebox{0.78\textwidth}{!}{%
  \begin{tikzpicture}[x=1.25cm,y=1.0cm]
    \draw[-Latex, thick] (0.8,0.8) -- (5.7,0.8) node[right, font=\scriptsize] {Probabilidad};
    \draw[-Latex, thick] (0.8,0.8) -- (0.8,5.55) node[above, font=\scriptsize] {Impacto};
    \foreach \x in {1,...,5} {
      \draw[linegray] (\x,0.8) -- (\x,5.3);
      \node[below, font=\scriptsize] at (\x,0.75) {\x};
    }
    \foreach \y in {1,...,5} {
      \draw[linegray] (0.8,\y) -- (5.4,\y);
      \node[left, font=\scriptsize] at (0.75,\y) {\y};
    }
    \fill[accentgreen!10] (0.8,0.8) rectangle (2.6,2.6);
    \fill[accentorange!12] (2.6,0.8) rectangle (4.2,4.2);
    \fill[accentred!12] (4.2,0.8) rectangle (5.4,5.3);
    \fill[accentred!12] (0.8,4.2) rectangle (5.4,5.3);
    \foreach \x/\y/\id in {4.6/4.7/R1,3.9/3.6/R2,3.2/3.2/R3,3.5/2.8/R4,2.8/3.9/R5,2.4/3.3/R6} {
      \node[circle, fill=uddblue, text=white, inner sep=2.4pt, font=\scriptsize\bfseries] at (\x,\y) {\id};
    }
    \node[font=\scriptsize, text=accentgreen] at (1.7,1.25) {Bajo};
    \node[font=\scriptsize, text=accentorange] at (3.35,2.15) {Medio};
    \node[font=\scriptsize, text=accentred] at (4.75,4.95) {Crítico};
  \end{tikzpicture}}
  \caption{Mapa de riesgos probabilidad-impacto para priorizar gestión del proyecto.}
\end{figure}

\section{Definición de Done por Hito}

\subsection*{M1: Verificación estructural}
El modelo contiene recursos, entidades, inventario, demanda preliminar y flujo sin errores de conexión. Se documentan supuestos iniciales y brechas abiertas.

\subsection*{M2: Validación cuantitativa inicial}
El escenario base reproduce aproximadamente los 8.200 kg/día, respeta patrones horarios y genera KPIs consistentes. La inconsistencia de capacidad de horno queda resuelta o explicitada.

\subsection*{M3: Recomendación preliminar}
La matriz de escenarios está ejecutada, el cuello de botella principal está identificado y se seleccionan tres configuraciones candidatas con evidencia.

\subsection*{M4: Calidad documental}
El informe técnico y académico contienen trazabilidad, anexos, tablas, supuestos, resultados, limitaciones y recomendación. La revisión cruzada no presenta no conformidades mayores.

\subsection*{M5: Entrega final}
El paquete final incluye modelo SIMIO ejecutable, informes, anexos de datos, bitácora de corridas, resultados y material de defensa.

\section{Checklist de Entrega Profesional}

\begin{enumerate}
  \item Modelo SIMIO ejecutable y versionado.
  \item Informe técnico orientado a gerencia.
  \item Informe académico con proceso, decisiones y reflexión crítica.
  \item Anexos con datos, parámetros, supuestos y distribuciones.
  \item Matriz de trazabilidad requisito--modelo--KPI--evidencia.
  \item Bitácora de corridas con versión, semillas, réplicas y resultados.
  \item Tablero de KPIs por escenario y por tipo de pan.
  \item Análisis de sensibilidad y cuello de botella.
  \item Recomendación final con frontera costo-servicio.
  \item Material de defensa: problema, evidencia, recomendación, riesgos e implementación.
\end{enumerate}

\section{Narrativa para Defensa}

La defensa debe comunicar el proyecto como una decisión operacional basada en evidencia:

\begin{quote}
La panadería enfrenta quiebres de stock porque la demanda es variable y la producción opera con recursos compartidos y restricciones de horneo. El modelo SIMIO permite reproducir el sistema, medir sus cuellos de botella y comparar configuraciones de recursos y políticas de producción. La recomendación final se sustenta en réplicas, intervalos de confianza, trazabilidad de datos y un balance explícito entre nivel de servicio, costo relativo y riesgo de sobreproducción.
\end{quote}

\section{Supuestos Iniciales}

\begin{itemize}
  \item No se consideran fallas de equipos.
  \item No hay sustitución entre tipos de pan cuando existe quiebre.
  \item El producto no se deteriora dentro del horizonte diario.
  \item Los hornos son equivalentes entre sí.
  \item La demanda del día base corresponde a 8.200 kg.
  \item Si no hay costos unitarios reales, se utiliza índice de costo relativo.
  \item La capacidad de horno se congelará solo después de resolver la inconsistencia entre enunciado y archivo de parámetros.
\end{itemize}

\section{Cierre}

Este roadmap establece una ruta de ejecución profesional para transformar el enunciado en un modelo de simulación útil, auditable y defendible. La calidad de la entrega dependerá de tres condiciones: consistencia de datos, disciplina de trazabilidad y validación estadística suficiente. Bajo esas condiciones, el equipo podrá entregar una recomendación operacional rigurosa y comprensible para una audiencia académica y gerencial.

\end{document}
